\section{Introduction}
\label{sec:intro}
In software development, tests are essential for identifying defects and verifying that fixes work correctly \cite{kapfhammer2004software}. Ideally, test results remain stable across runs of an unchanged codebase. However, flaky tests, that is, tests that pass or fail nondeterministically on the same version of the code, are a common reality \cite{parry2021survey, pinto2020vocabulary}. Their causes are diverse, including concurrency issues, asynchronous operations, order dependencies between tests, improper management of external resources, dependence on the system clock, and the use of random values \cite{luo2014empirical}, as well as overly narrow ranges of expected values and timeout-related issues \cite{eck2019understanding}.

Although flaky tests do not directly affect software functionality, they impose a substantial cost on delivery pipelines. Developers report wasted effort analyzing failures, ignored failures, and loss of trust in the test suite \cite{eck2019understanding, gruber2022survey}, and at Google about 16\% of tests show some level of flakiness, with 2\% to 16\% of compute resources spent on rerunning them \cite{micco2017state}. Detecting flakiness through reruns is expensive, since tests that fail rarely require many executions \cite{lam2020study}. Static prediction from the test code is therefore attractive: it requires running nothing, and its features can be inspected by developers.

Pinto et al.~\cite{pinto2020vocabulary} showed that flaky tests share a characteristic vocabulary, that is, identifiers and keywords that occur disproportionately in them, and achieved an F-score of 0.95 in Java. The idea was replicated for Java with more classifiers and an inter-project evaluation \cite{camara2021revisiting}, for Python \cite{haben2021replication}, JavaScript \cite{soratto2023vocabulary}, and Swift \cite{medeiros2026vocabulary}, and extended to five languages \cite{ahmad2025understanding}. These studies also expose its limits. Performance drops sharply across projects \cite{camara2021revisiting, soratto2023vocabulary} and over time \cite{haben2021replication}, the predictions were not actionable in an industrial setting \cite{berndt2023vocabulary}, and code-based detectors do not beat a majority baseline once projects are kept apart \cite{gultekin2026limits}.

To the best of our knowledge, vocabulary-based prediction has never been evaluated for C\#, even though flakiness in .NET projects is well documented \cite{lam2019root, lam2020study, leesatapornwongsa2022flakerepro} and the language has relevant particularities. Its concurrency model centers on the Task-based asynchronous pattern (\texttt{Task}, \texttt{async}/\texttt{await}, \texttt{CancellationToken}) and the thread pool, which introduces its own timing and scheduling dependencies. In tests from Microsoft .NET projects, concurrency is the leading cause of flakiness, with fixed waits and timeouts as a frequent pattern \cite{leesatapornwongsa2022flakerepro}. Test suites are written with several testing frameworks (xUnit, NUnit, and MSTest) and run on two test platforms (VSTest and the newer Microsoft.Testing.Platform), and libraries are often compiled for multiple target frameworks, so the same test may run under different runtimes. \Den{Nenhum dos 17 projetos usa MSTest (Seção~\ref{sec:subjects}), e as particularidades deste parágrafo (frameworks, multi-target, locale/time zone, containers) não são analisadas depois. Ou se faz uma análise por framework/TFM (os tokens e o desempenho mudam entre xUnit e NUnit?), ou se enxuga o parágrafo para não prometer o que o artigo não entrega.} Integration suites also depend on containerized services and on environment behaviors such as locale, time zone, and file system paths. If the lexical patterns that signal flakiness in Java, JavaScript, and Swift do not transfer to C\#, vocabulary-based tools built around them would offer little help to .NET developers. This is the gap this work addresses. \Den{A lacuna é formulada como ``os padrões transferem para C\#?'', mas nenhuma RQ testa transferência. Sugestão de RQ nova: comparar sistematicamente o vocabulário de C\# com as listas publicadas para Java (Pinto), JS (Soratto), Swift (Medeiros) e Python (Haben), por exemplo pela sobreposição dos top-$k$ tokens após normalização (sleep, thread, wait, async, timeout...). Uma versão mais forte: treinar com dados de outra linguagem e testar em C\#. Isso também justificaria o título.}

To address it, we investigate static prediction of flaky tests in C\# from the vocabulary of test source code. Our approach, organized in six steps (Section~\ref{sec:approach}), selects open-source C\# projects and systematically reruns their test suites, 100 times in an automated fashion and again 20 times under controlled isolation; labels each test as flaky, broken, or stable according to its outcomes; resolves inherited tests to the method that declares them; extracts the vocabulary and structural features of each test with the Roslyn compiler platform; and trains classifiers on this vocabulary to predict a test's propensity to be flaky without executing it, while also extracting an interpretable ranking of the tokens most strongly associated with flakiness. We ask the following research questions:

\begin{description}[leftmargin=1.2em,labelindent=0pt,itemsep=1pt]
  \item[RQ1.1] How accurately can flaky tests in C\# be predicted from the vocabulary of their source code?
  \item[RQ1.2] Do the labels and results remain stable between the 100-run and the 20-run rounds?
  \item[RQ2] Does vocabulary-based prediction outperform the majority, random, vocabulary-threshold, and structure-only baselines?
  \item[RQ3] Which tokens are most strongly associated with flakiness?
  \item[RQ4] When and why does the classifier fail?
\end{description}

We applied this methodology to 17 projects, with 19,211 test cases and 12,665 labeled test methods, of which 30 are flaky and 12,635 are stable (0.24\%). Among the eight classifiers evaluated, tree ensembles performed best under cross-validation, with Extra Trees reaching an MCC of 0.53 and an AUC of 0.86, and Random Forest an MCC of 0.51 and an AUC of 0.87. These results are well above random guessing (MCC 0.02), but only moderately above a structure-only baseline (MCC 0.36) and a vocabulary threshold rule (MCC 0.40) (RQ1.1, RQ2). With the labels of the 20-run round, the structure-only model catches up with the ensembles, so the advantage of vocabulary depends on the round that produced the labels (RQ1.2). When each repository is held out and evaluated at its natural base rate, AUC drops to 0.75--0.78 and precision stays at 2--3\%. The approach is therefore better suited to ranking tests and prioritizing reruns than to gating a pipeline, since the top quartile of each repository's ranking contains 23 of the 30 flaky tests. The signal rests on Task-based concurrency, time handling, and resource lifecycles (RQ3), and the model fails when the cause lies outside the test body or when asynchronous constructs are used deterministically (RQ4).

The main contributions of this work are:
\begin{packed_itemize}
  \item[$\star$] A dataset of flaky, broken, and stable C\# tests, built from 100 automated and 20 controlled reruns of 17 open-source projects, with inherited tests resolved to the methods that declare them.
  \item[$\star$] An evaluation of eight classifiers for vocabulary-based static prediction in C\#, under repeated cross-validation and under cross-repository validation at the natural base rate, showing that a lexical signal exists in C\# but is modest under realistic evaluation.
  \item[$\star$] A C\# flakiness vocabulary, which separates frequent general markers from rare, specific ones and shows that no robust marker of stability emerges.
  \item[$\star$] An analysis of the model's errors in the test code, with sensitivity analyses without attributes and with helper-method expansion, which outline the limits of lexical prediction.
  \item[$\star$] An assessment of label robustness, comparing the two rerun rounds (which agree on only 13 of the 36 flaky tests), and documentation of the pitfalls of rerun studies in .NET that distorted our initial counts (Table~\ref{tab:armadilhas}).
\end{packed_itemize}
\Den{Posicionamento para o ISSTA: replicar o método numa linguagem nova tende a ser visto como incremental (perfil mais de MSR/ESEM/SAST). As contribuições mais originais aqui são (i) o protocolo e as armadilhas de reruns em .NET, (ii) a sensibilidade dos rótulos à rodada e (iii) a avaliação realista (LOPO, taxa natural, ranking). Sugiro organizar a narrativa em torno disso e deixar ``vocabulário em C\#'' como veículo. A última contribuição enquadra a sobreposição 13/36 como ``robustez''; ver o comentário na RQ1.2.}
